\reading{IM-12}{Performing Due Diligence on Specific Managers and Funds}{STRIKE FIRST}{3}

\srcnote{Kevin R. Mirabile, \textit{Hedge Fund Investing: A Practical Approach to
Understanding Investor Motivation, Manager Profits, and Fund Performance}, 2nd
Edition (Hoboken, NJ: Wiley Finance, 2016), Chapter 12. Printed as Chapter 10 in
the GARP source volume --- the last GARP-primary reading in this area.}

\begin{imkeybox}[title={The seven objectives in this reading}]
\begin{enumerate}[label=\textbf{\alph*.}]
\item Identify reasons for the failures of hedge funds in the past.
\item Explain elements of the due diligence process used to assess investment managers.
\item Identify themes and questions investors can consider when evaluating a hedge fund manager.
\item Describe criteria that can be evaluated in assessing a hedge fund's risk management process.
\item Explain how due diligence can be performed on a hedge fund's operational environment.
\item Explain how a hedge fund's business model risk and its fraud risk can be assessed.
\item Describe elements that can be included as part of a due diligence questionnaire.
\end{enumerate}
\end{imkeybox}

\begin{imtrapbox}[breakable=false,title={Two layers, two wrong reading numbers, and a wrong area}]
The crash layer labels this \emph{Reading 92} --- which is IM-10, private credit,
covered correctly as 92 elsewhere in the same document --- and runs its page
header as \term{``Credit Risk Measurement \& Management''}, an entirely different
exam area. The lecture layer labels the same material \emph{LO 87.a--f}, and 87 is
IM-5. Neither layer is right, and the two are not wrong in the same way. The
crash layer is also the only one of the two that reaches objective \term{g}.
\end{imtrapbox}

% ------------------------------------------------------------------
\basics{Basics: what due diligence is}

Due diligence is the process of evaluation and analysis an investor follows to get
comfortable with \term{a strategy, a manager, and a fund} before investing.
Investors would be wise to assess \term{all three skill sets} needed to run a
modern hedge fund --- \term{investment, operational and business} skills ---
before committing capital. It includes learning why and how a fund came into
being, the skills its founders claim to have mastered, the timeliness, accuracy
and consistency of manager and fund information, and the reliability and
independence of service providers.

\begin{imkeybox}
The three skill sets are not decoration: they map directly onto the reading's
structure. \term{Investment} skill is IM-12 c and d; \term{operational} skill is
IM-12 e; \term{business} skill is IM-12 f.
\end{imkeybox}

% ------------------------------------------------------------------
\lo{IM-12 a}{Identify reasons for the failures of hedge funds in the past.}

\begin{enumerate}[label=\alph*)]
\item Bad investment decisions.
\item All sorts of \term{frauds} --- accounting frauds, valuation frauds, or
misappropriation of funds.
\item \term{Excessive leverage}, improbable probabilities, unexpected events, and
tail risk.
\item Getting caught in \term{squeezes} by other hedge funds.
\item Lack of supervision or compliance controls related to \term{insider
trading}.
\item A flood of \term{unanticipated withdrawals} of capital at the least
opportune time; funds can fail when liquidity dries up and they cannot meet
redemptions.
\item Funds can fail because of \term{their own actions or the acts of others}.
\end{enumerate}

\begin{imtrapbox}[breakable=false,title={Only the first is about being wrong}]
Of the seven reasons, exactly one --- bad investment decisions --- is about
getting the trade wrong. The other six are about fraud, leverage, liquidity,
supervision and other people. That ratio is the argument for the rest of the
reading: due diligence on the investment process alone examines the least
frequent cause of failure.
\end{imtrapbox}

% ------------------------------------------------------------------
\lo{IM-12 b}{Explain elements of the due diligence process used to assess investment managers.}

\renewcommand{\arraystretch}{1.2}
\noindent\begin{tabularx}{\textwidth}{@{}>{\raggedright\arraybackslash}p{3.0cm}X@{}}
\toprule
\textbf{\sffamily\small Element} & \textbf{\sffamily\small What is asked} \\
\midrule
Investment process & What is your strategy, and how does it work? How is equity ownership allocated among the portfolio management, trading and research teams? \term{Is the track record reliable?} Who are the principals, and are they trustworthy? \\
Related risk controls & How is risk measured and managed? \term{How are securities valued?} What is the portfolio leverage and liquidity? Does the strategy expose the investor to tail risk? How often do investors get risk reports, and what do they include? Do the fund terms make sense for the strategy? \\
Operational environment & Internal control assessment; documents and disclosures; service provider evaluation. \\
Model risk and fraud risk & Business model risk deals with the risks of simply \emph{running the business} of being a hedge fund. Fraud risk persists \term{even where due diligence on the investment, risk management and operational practices is complete}, and even where the business model is fully understood. \\
\bottomrule
\end{tabularx}
\renewcommand{\arraystretch}{1}
\figcap{The crash layer compresses these four into three --- investment background
and performance; risk controls and processes; business model and operations ---
which folds operational due diligence into the third. The lecture layer's
four-way split is the fuller one and is used here. Note the last row: the source
is explicit that thorough due diligence \emph{does not eliminate} fraud risk.}

% ------------------------------------------------------------------
\lo{IM-12 c}{Identify themes and questions investors can consider when evaluating a hedge fund manager.}

Manager evaluation breaks into \term{four areas: strategy, ownership, track
record, and investment management.}

\begin{itemize}
\item \term{Strategy.} Any particular investment strategy or current investment
trend; investment style; extent of turnover and liquidity; mechanism to limit
potential loss; research and analysis done by the investment team; any use of
derivatives; the trade execution process; extent of investment in private
companies; the trade-off between current returns and long-term fund growth; and
\term{has the fund ever been shut down}.
\item \term{Ownership.} Ownership interests help \term{align the interests} of the
investment team and investors, and help attract and retain quality staff.
Investors should ask whether traders, portfolio managers and research analysts
have ownership interests in the firm.
\item \term{Track record.} Comparison with peers; \term{past performance audited
by a third party}; trend and attribution analysis; returns relative to the size of
the investment; and \term{whether the employees who generated the past returns are
still with the fund}.
\item \term{Investment management.} The manager interview (strategy, and how the
manager coped in tough market conditions); \term{reference checks} with former
employers, current and former colleagues, and current and former investors; and
\term{background checks} --- a background report, review of the \term{ADV form},
past litigation or criminal behaviour, the manager's personal credit and
bankruptcy report, enquiries with past colleagues, and involvement in related
party transactions.
\end{itemize}

\begin{imtrapbox}[breakable=false,title={The personal credit report is not an accident}]
Background checks reach into the manager's \emph{personal} credit and bankruptcy
history, not just the firm's. A manager under personal financial pressure is a
fraud risk factor, which is why this sits in IM-12 c and reappears in IM-12 f.
\end{imtrapbox}

% ------------------------------------------------------------------
\lo{IM-12 d}{Describe criteria that can be evaluated in assessing a hedge fund's risk management process.}

\begin{enumerate}[label=\alph*)]
\item \term{Risk.} Assess all systematic and unsystematic risks; are there written
policies; is there a risk committee; the extent of risk management culture; the
information technology used to quantify risk; the existence and structure of any
risk model.
\item \term{Security valuation.} Identify the \term{proportion of assets whose
market price is observable}; examine the \term{independence} of valuation;
determine whether price is \term{overridden} for valuation purposes.
\item \term{Portfolio leverage and liquidity.} Sources of current and future
leverage; current liquidity level; adjustment to return for using excessive
leverage.
\item \term{Exposure to tail risks.} Does the return distribution show skewness or
kurtosis; preparedness to meet a tail event.
\item \term{Risk reports.} Review key reports before investing; analyse key
metrics and benchmark them against peers.
\item \term{Consistency of fund terms with the investment strategy.} Examine the
fee structure; identify any additional fee above a threshold; is a high fee being
paid \term{for creating alpha or beta}; identify limitations on redemption.
\end{enumerate}

\begin{imkeybox}[title={The two questions that do the most work}]
\term{What proportion of assets has an observable market price?} and
\term{is the price ever overridden?} Together they determine whether the reported
NAV is a measurement or an opinion --- which is the same issue as the stale
valuations of IM-10 c and the DPI-against-TVPI test of IM-11 d.
\end{imkeybox}

% ------------------------------------------------------------------
\lo{IM-12 e}{Explain how due diligence can be performed on a hedge fund's operational environment.}

Three focus areas: \term{internal control assessment, documents and disclosure,
and service provider evaluation}.

\renewcommand{\arraystretch}{1.2}
\noindent\begin{tabularx}{\textwidth}{@{}>{\raggedright\arraybackslash}p{2.8cm}X@{}}
\toprule
Internal control assessment & \term{Personnel:} qualifications, attitudes and training of key personnel including the CEO, and their commitment to controls and compliance. \term{Policies and procedures:} review policies, procedures and audit reports for control effectiveness, and ensure they are actually followed. \term{Compliance system:} the code of ethics and restrictions on trading and related-party transactions. \term{Counterparty risk:} how the fund diversifies counterparties and monitors them. \term{Corporate governance:} its effectiveness, and how breaches are addressed and prevented. \\
Documents and disclosure & \term{Document verification:} confirm legal counsel's involvement in document creation; \emph{physically inspect} documents for unauthorized changes; verify terms across documents. \term{Conflicts and risks:} scrutinise disclosed conflicts; review risk disclosures; assess the manager's authority, limitations and reporting duties. \term{Financial statements:} ensure \term{unqualified audit opinions}; check the balance sheet and income statement align with the strategy; examine footnotes; \term{validate manager fees, especially in loss years}; investigate general partner contributions and fund withdrawals. \\
Service provider evaluation & Examine internal control letters and audited financial statements of third-party providers, and conduct \term{in-person discussions} about their role and performance. \\
\bottomrule
\end{tabularx}
\renewcommand{\arraystretch}{1}
\figcap{Three tests that are hard to fake, and one that is not. Physical
inspection of documents, in-person discussions with service providers and the
validation of fees in \emph{loss} years all require someone to be in a room. A
review of stated policies does not --- which is why the source pairs it with
``ensure they're followed and effective.''}

% ------------------------------------------------------------------
\lo{IM-12 f}{Explain how a hedge fund's business model risk and its fraud risk can be assessed.}

\begin{imdefbox}[title={Business model risk}]
The risk associated with a fund's ability to \emph{operate successfully and
generate returns} --- as a business, not as a portfolio. Typical risks:
\term{insufficient cash}, \term{lack of a succession plan}, and \term{excessive
redemptions in a short time}.
\end{imdefbox}

\begin{itemize}
\item \textbf{Assessing business model risk:} examine revenue and expense
stability; calculate \term{the reliance on variable performance fees}; evaluate
the gap between management fees and expenses; check working capital adequacy;
assess budgeting practices and \term{breakeven points}; consider capacity to
handle increased assets; verify \term{key person insurance} and succession
planning.
\item \textbf{Assessing fraud risk:} frequent \term{related-party transactions};
high illiquidity, \emph{especially with manager-only valuations}; frequent
litigation, particularly fraud claims; \term{unreasonably high stated returns};
the manager's \term{personal trading} in fund-related securities; and frequent
shorting transactions.
\item \textbf{Mitigating fraud risk:} check regulatory infractions on the SEC
website; search court records for past litigation and bankruptcy history; verify
the \term{competence and independence of service providers}; and conduct thorough
background checks on the manager.
\end{itemize}

\begin{imtrapbox}[breakable=false,title={Sibling pair: the two risks share indicators}]
Illiquidity appears in both lists, and so does the manager's behaviour. The
discriminator is the \emph{question being asked}. Business model risk asks
\term{can this firm survive as a business}; fraud risk asks \term{is someone
lying}. A fund with unstable revenue and no succession plan has business model
risk; a fund whose illiquid assets are valued by the manager alone has fraud risk.
Madoff, in IM-14, is the case where the second was ignored.
\end{imtrapbox}

% ------------------------------------------------------------------
\lo{IM-12 g}{Describe elements that can be included as part of a due diligence questionnaire.}

\srcnote{Covered by the crash layer only. The lecture layer stops at objective f.}

\begin{enumerate}[label=\alph*)]
\item \term{General information.} Manager registration, ownership structure, key
personnel, past performance, and contact details.
\item \term{Fund information.} Fees, lockup periods, redemption rules, service
providers including auditors and legal counsel, historical performance, and asset
details.
\item \term{Execution and trading.} Transaction processing speed and accuracy, and
\term{potential conflicts with service providers}.
\item \term{Research policy.} The fund's sources of research information.
\item \term{Compliance.} The existence of legal counsel and anti-money laundering
policies.
\end{enumerate}

\begin{imtrapbox}[title={Consolidated trap summary --- IM-12}]
\begin{itemize}
\item \term{Role.} The three skill sets are \emph{investment, operational and
business}. Substituting ``financial'' or ``compliance'' for any of the three is
the standard corruption.
\item \term{Scope, IM-12 a.} Six of the seven failure causes are not about
investment decisions.
\item \term{Definition, IM-12 b.} Fraud risk survives complete due diligence on
investment, risk management and operational practice. The source says so
explicitly.
\item \term{Role, IM-12 c.} Track record questions include whether the
\emph{people} who generated the returns are still there, and whether performance
was audited by a \emph{third party}.
\item \term{Definition, IM-12 d.} The two valuation questions are what
\emph{proportion} of assets has an observable price, and whether price is ever
\emph{overridden}.
\item \term{Sibling, IM-12 f.} Business model risk is about surviving as a
business; fraud risk is about someone lying. Illiquidity is an indicator of both.
\item \term{Scope, IM-12 g.} The questionnaire is in the crash layer only and is
the objective most often left uncovered.
\end{itemize}
\end{imtrapbox}

% ==================================================================
